% =============================================================================
% Project : THE_PHILOSOPHY_OF_VIABILITY
% Book    : Book I — Ontology of Viability
% Book ID : VS.B01
%
% File    : VS.B01_ONTOLOGY_OF_VIABILITY/2_MAINMATTER/PART_I_IDENTITY_AND_MANIFESTATION/CHAPTER_07_FORMAL_AND_LOGICAL_RECONSTRUCTION/sections/section_14_formal_reconstruction_of_expression_and_manifestation.tex
%
% Section / Chapter:
%   Chapter 07, Section 14
%   Formal Reconstruction of Expression and Manifestation
%
% Purpose :
%   To reconstruct Expression and Manifestation as distinguishable
%   ontological relations, specify their formal roles and possible
%   connections, and prevent their reduction to one another or to
%   identity, determination, or cognition.
%
% =============================================================================

\section{Formal Reconstruction of Expression and Manifestation}
\label{sec:formal-reconstruction-of-expression-and-manifestation}

\subsection{The Problem of Formal Distinction}

Expression and Manifestation concern how an existent
is articulated or becomes manifest in a specified
ontological context. Both concepts are necessary
for examining the ways in which an existent may
be present, articulated, disclosed, or encountered
under particular conditions.

Their proximity, however, creates a conceptual
and formal difficulty.

If Expression and Manifestation are treated as
interchangeable, the ontology loses the ability
to distinguish different relations through which
an existent may be articulated or manifested.

If they are treated as wholly unrelated, the
formal system becomes unable to represent
their possible connections.

The required reconstruction must therefore
preserve two conditions simultaneously:

\begin{enumerate}
    \item Expression and Manifestation must
    remain analytically distinguishable.
    \item Their possible ontological relations
    must remain formally examinable.
\end{enumerate}

The central thesis of this section is that
Expression and Manifestation require distinct
formal representations, while their relationship
must be determined through specified conditions
rather than assumed through terminological
similarity.

This thesis does not establish that every
existent must be expressed or manifested in
the same manner. Nor does it assume that
Expression and Manifestation are necessarily
coextensive.

The purpose of formalization is to represent
the relevant relations precisely enough to
examine their conditions, dependencies, and
limits.

\subsection{The Formal Domains of Expression and Manifestation}

Let:

\[
    D_E
\]

denote a specified domain of existents.

Let:

\[
    D_X
\]

denote a formal domain of expression instances,
and let:

\[
    D_M
\]

denote a formal domain of manifestation instances.

The symbols \(D_X\) and \(D_M\) designate
formal domains for the purpose of analysis.
Their introduction does not establish that
expressions and manifestations are independent
substances or entities.

They may represent events, configurations,
relations, modes, occurrences, or other
formally distinguishable items, depending
on the interpretation.

The formal system must specify which
interpretation is intended in each application.

Let:

\[
    x \in D_E,
    \qquad
    e \in D_X,
    \qquad
    m \in D_M.
\]

The variables represent, respectively,
an existent, an expression instance, and
a manifestation instance.

The use of separate domains prevents
the formal language from treating the
expression and manifestation variables
as interchangeable merely because
they may be associated with the same
existent.

A many-sorted language is appropriate
where the distinctions among these
formal categories are necessary for
valid inference.

A single-sorted language may also be
used if its predicates and quantifiers
preserve the required distinctions.

The choice of formal language does not
settle the ontological status of the
categories. It determines how their
relations are represented.

\subsection{Expression as a Distinct Relation}

Let:

\[
    \operatorname{Expr}(x,e)
\]

represent a specified expression relation
between existent \(x\) and expression
instance \(e\).

The expression relation must be interpreted
according to the ontological account
adopted for Expression.

The formal symbol alone does not determine
whether Expression concerns articulation,
configuration, articulation of a determination,
or another specified relation.

Its meaning must be fixed through the
definitions and arguments of the ontology.

The formula:

\[
    \operatorname{Expr}(x,e)
\]

therefore expresses a relation whose
conditions are to be specified, rather
than a complete definition of Expression.

At least three elements may require
distinction in a particular expression
claim:

\begin{itemize}
    \item the existent to which the
    expression is related;
    \item the expression instance under
    consideration;
    \item the respect or condition under
    which the expression relation holds.
\end{itemize}

A more refined notation may be introduced:

\[
    \operatorname{Expr}(x,e,\rho),
\]

where \(\rho\) identifies a specified
respect of expression.

The parameter \(\rho\) is not an
additional ontological entity. It is
a formal device for distinguishing
different interpretations of the
expression relation.

For example, two expression claims
concerning the same existent may
differ in their respective conditions
without thereby establishing two
different existents or two different
identities.

The formal system must therefore
avoid treating every difference in
expression as a difference in identity.

\subsection{Manifestation as a Distinct Relation}

Let:

\[
    \operatorname{Man}(x,m)
\]

represent a specified manifestation
relation between existent \(x\) and
manifestation instance \(m\).

The relation concerns the manifestation
of an existent under the interpretation
adopted for the relevant ontological
analysis.

The formula does not, by itself,
determine the conditions under which
a manifestation occurs, the nature
of what is manifested, or the relation
between manifestation and the
existent's identity.

Those matters require further
specification.

A respect-sensitive form may be written:

\[
    \operatorname{Man}(x,m,\sigma),
\]

where \(\sigma\) identifies a specified
respect or condition of manifestation.

The parameter \(\sigma\) may distinguish
different manifestation conditions
without presupposing that those
conditions constitute different
identities.

The formal reconstruction must
also distinguish manifestation
instances from the objects, contexts,
or conditions with which they may
be associated.

A manifestation relation is not
automatically a relation of
representation, perception, or
description.

Such interpretations may be relevant
in particular cases, but they cannot
be imposed as universal conditions
of manifestation.

The ontology under consideration
does not make the existence of an
existent or its identity dependent
upon a knower's recognition of
its manifestation.

The formal language must therefore
remain capable of representing
manifestation without introducing
cognition as a necessary condition.

\subsection{The Non-Identity of Expression and Manifestation}

The formal reconstruction requires
distinct predicates for Expression
and Manifestation:

\[
    \operatorname{Expr}(x,e)
\]

and:

\[
    \operatorname{Man}(x,m).
\]

The use of distinct predicates is
not merely a stylistic choice.
It preserves the possibility that
the two relations have different
conditions and different formal
consequences.

The formal system must not adopt
the unrestricted equivalence:

\[
    \operatorname{Expr}(x,e)
    \leftrightarrow
    \operatorname{Man}(x,e).
\]

Apart from the possible type mismatch
between \(D_X\) and \(D_M\), this
biconditional would assume that
Expression and Manifestation have
the same extension and conditions.

Such an assumption requires an
independent argument and cannot
be derived from their conceptual
proximity.

Nor should the formal system
assume that every expression
necessarily produces a manifestation:

\[
    \operatorname{Expr}(x,e)
    \not\Rightarrow
    \exists m\,
    \operatorname{Man}(x,m).
\]

Likewise, it should not assume
that every manifestation necessarily
constitutes an expression:

\[
    \operatorname{Man}(x,m)
    \not\Rightarrow
    \exists e\,
    \operatorname{Expr}(x,e).
\]

These formulas express restrictions
on unrestricted inference. They
do not assert that Expression
and Manifestation never occur
together.

Their possible connection remains
open to examination through
additional relations and principles.

\subsection{The Possible Relation Between Expression and Manifestation}

The distinction between Expression
and Manifestation does not prevent
their connection.

Let:

\[
    \operatorname{Rel}_{XM}(e,m)
\]

represent a specified relation
between expression instance \(e\)
and manifestation instance \(m\).

The relation may be interpreted
according to the particular
ontological account under
consideration.

Its introduction allows the
formal system to examine
whether a manifestation is
associated with an expression,
whether an expression is
manifested under specified
conditions, or whether the
two relations occur together
without a direct dependence.

A more specific relation may
be written:

\[
    \operatorname{ManifestsThrough}(x,e,m).
\]

This formula may represent
a specified case in which
an existent is manifested
through an expression instance.

It is not a universal definition
of Manifestation or Expression.

Its validity depends on the
conditions assigned to the
relation.

The distinction between the
general relations and their
possible connection is important.

The general expression relation
may be represented as:

\[
    \operatorname{Expr}(x,e),
\]

and the general manifestation
relation as:

\[
    \operatorname{Man}(x,m).
\]

Their connection may be represented
separately as:

\[
    \operatorname{Rel}_{XM}(e,m).
\]

This structure avoids defining
one relation through the other.

It also allows different models
to represent different forms
of connection.

For example, a model may
interpret the two relations
as systematically associated
under certain conditions.

Another model may interpret
them as independently satisfiable
relations.

The formal framework permits
both possibilities until an
ontological argument establishes
a more determinate account.

\subsection{Expression, Manifestation, and Determination}

Expression and Manifestation
must also be distinguished
from Determination.

Let:

\[
    \operatorname{Det}(x,d)
\]

represent a specified determination
relation between existent \(x\)
and determination instance or
formal determination reference \(d\).

The notation is provisional
and must be interpreted in
accordance with the ontology's
account of Determination.

The formal system must not
assume:

\[
    \operatorname{Expr}(x,e)
    \leftrightarrow
    \operatorname{Det}(x,e),
\]

or:

\[
    \operatorname{Man}(x,m)
    \leftrightarrow
    \operatorname{Det}(x,m).
\]

These equivalences would collapse
different ontological questions
into one formal relation.

A determination claim concerns
a specified determination of
an existent.

An expression claim concerns
the relation represented by
Expression.

A manifestation claim concerns
the relation represented by
Manifestation.

The exact connections among
these notions require independent
analysis.

It may be possible for a
particular determination to
be expressed or manifested
under specified conditions.
However, such a possibility
does not establish that
Determination is identical
to Expression or Manifestation.

The formal system may therefore
introduce bridge relations:

\[
    \operatorname{ExprDet}(d,e),
\]

\[
    \operatorname{ManDet}(d,m),
\]

where these relations are
required by a particular
interpretation.

Their presence must be justified
rather than presumed.

\subsection{Expression, Manifestation, and Identity}

The relation between Expression,
Manifestation, and Identity
requires particular care.

The foundational framework
maintains the analytical
distinction and ontological
inseparability of Being and
Identity, together with the
absence of ontological priority
between them.

Expression and Manifestation
must be formalized consistently
with these commitments.

Let:

\[
    \operatorname{Id}(x,i)
\]

represent a specified formal
relation between existent \(x\)
and identity reference \(i\).

This notation is introduced
for formal examination. It
does not construe identity
as an additional property
attached to an otherwise
complete existent.

The formal system must not
identify Expression with
Identity:

\[
    \operatorname{Expr}
    \not\equiv
    \operatorname{Id}.
\]

Nor may it identify
Manifestation with Identity:

\[
    \operatorname{Man}
    \not\equiv
    \operatorname{Id}.
\]

These distinctions preserve
the different analytical
roles of the concepts.

An expression may be associated
with an existent whose identity
is under examination.

A manifestation may be associated
with the same existent.

Neither association establishes
that the expression or
manifestation constitutes
the identity itself.

Similarly, the presence of
different expressions or
manifestations does not
automatically establish
identity plurality.

The following implications
must therefore not be adopted
without further premises:

\[
    \operatorname{Pl}_{X}(x)
    \rightarrow
    \operatorname{Pl}_{I}(x),
\]

\[
    \operatorname{Pl}_{M}(x)
    \rightarrow
    \operatorname{Pl}_{I}(x).
\]

Here, \(\operatorname{Pl}_{X}\)
and \(\operatorname{Pl}_{M}\)
represent specified plurality
conditions concerning expression
and manifestation instances,
while \(\operatorname{Pl}_{I}\)
represents identity plurality.

The formal reconstruction
preserves the possibility
that several expressions
or manifestations may be
associated with an existent
without deciding the number
of identities involved.

The ontological question
of identity plurality remains
independent and must be
examined through its own
arguments.

\subsection{Multiplicity of Expression and Manifestation}

Expression and Manifestation
may each involve plurality.

Let:

\[
    \operatorname{Pl}_{X}(x)
\]

represent plurality of expression
instances associated with
existent \(x\).

A simple formal schema is:

\[
    \exists e_1,e_2 \in D_X
    \bigl(
        \operatorname{Expr}(x,e_1)
        \land
        \operatorname{Expr}(x,e_2)
        \land
        e_1 \neq e_2
    \bigr).
\]

Likewise, manifestation plurality
may be represented as:

\[
    \operatorname{Pl}_{M}(x)
\]

with the schema:

\[
    \exists m_1,m_2 \in D_M
    \bigl(
        \operatorname{Man}(x,m_1)
        \land
        \operatorname{Man}(x,m_2)
        \land
        m_1 \neq m_2
    \bigr).
\]

These formulas establish
formal plurality under the
chosen interpretation.

They do not determine
whether the instances differ
in kind, degree, configuration,
context, temporal occurrence,
or another respect.

Those distinctions must
be specified where they
matter to the ontological
argument.

Nor do the formulas establish
that all expression or
manifestation instances
are mutually independent.

A more refined formalization
may introduce relations
of dependence, ordering,
compatibility, or transformation.

For example:

\[
    \operatorname{DependsOn}(e_1,e_2)
\]

may represent a specified
dependence between expression
instances.

Similarly:

\[
    \operatorname{Compatible}(m_1,m_2)
\]

may represent a specified
compatibility relation between
manifestation instances.

These are optional formal
relations. They must not
be included as universal
assumptions without
ontological justification.

The central result is that
plurality of expression
and plurality of manifestation
must be examined independently
before their relation to
other forms of plurality
is established.

\subsection{Conditions, Contexts, and Respect}

Expression and Manifestation
may require the specification
of conditions under which
their respective relations
hold.

Let:

\[
    \kappa
\]

denote a specified condition
or context.

The expression relation
may be refined as:

\[
    \operatorname{Expr}(x,e,\kappa),
\]

and the manifestation relation
as:

\[
    \operatorname{Man}(x,m,\kappa).
\]

The parameter \(\kappa\)
is a formal placeholder
whose meaning must be
specified in each application.

It may refer to a relevant
ontological condition,
configuration, relation,
or context.

It must not automatically
be interpreted as a cognitive
or linguistic context.

The ontology must distinguish
conditions belonging to
the existent or its relations
from conditions introduced
only by an observer's
description.

The same existent may
therefore be examined
under different expression
or manifestation conditions
without assuming that every
change in condition constitutes
a change in identity.

Such a conclusion would
require an additional
argument concerning the
relation between the
relevant transformation
and identity.

\subsection{Formal Representation of Expression and Manifestation Structures}

The basic formal architecture
may be represented by
the following relations:

\[
    \operatorname{Expr}(x,e),
\]

\[
    \operatorname{Man}(x,m),
\]

\[
    \operatorname{Rel}_{XM}(e,m),
\]

\[
    \operatorname{Det}(x,d),
\]

\[
    \operatorname{Id}(x,i).
\]

The architecture distinguishes
five formal roles:

\begin{enumerate}
    \item the existent under
    examination;
    \item the expression relation;
    \item the manifestation relation;
    \item the possible relation
    between expression and
    manifestation;
    \item the independent
    determination and identity
    relations.
\end{enumerate}

The architecture is not
a causal diagram or a
hierarchy of ontological
priority.

The order in which the
relations are written
does not establish that
one relation produces
another.

Nor does the architecture
assume that every existent
must satisfy every relation
in every interpretation.

Its function is to preserve
the distinctions required
for valid formal analysis.

A more detailed system
may introduce additional
relations when the
ontological argument
requires them.

However, every additional
relation must have a
specified interpretation
and must not merely
rename an existing concept.

\subsection{Definitions and Propositions}

The following definitions
establish provisional
formal conventions for
the present reconstruction.

\paragraph{Definition 14.1: Expression relation.}

\(\operatorname{Expr}(x,e)\)
represents a specified
relation between existent
\(x\) and expression instance
\(e\), under the interpretation
adopted for Expression.

\paragraph{Definition 14.2: Manifestation relation.}

\(\operatorname{Man}(x,m)\)
represents a specified
relation between existent
\(x\) and manifestation
instance \(m\), under
the interpretation adopted
for Manifestation.

\paragraph{Definition 14.3: Expression--Manifestation relation.}

\(\operatorname{Rel}_{XM}(e,m)\)
represents a specified
relation between an
expression instance and
a manifestation instance.

\paragraph{Definition 14.4: Expression plurality.}

\(\operatorname{Pl}_{X}(x)\)
represents the existence
of at least two formally
distinguishable expression
instances associated with
\(x\), under the specified
expression relation.

\paragraph{Definition 14.5: Manifestation plurality.}

\(\operatorname{Pl}_{M}(x)\)
represents the existence
of at least two formally
distinguishable manifestation
instances associated with
\(x\), under the specified
manifestation relation.

These definitions specify
the formal vocabulary.
They do not establish
the ontological truth
of every proposition
expressed through it.

The following propositions
identify the principal
results and constraints
of the reconstruction.

\paragraph{Proposition 14.1: Expression and Manifestation require distinct formal representations.}

A formal reconstruction
that uses one undifferentiated
relation for Expression
and Manifestation cannot
represent their possible
differences without
additional structure.

\paragraph{Proposition 14.2: Expression does not automatically entail Manifestation.}

The expression relation
alone does not establish
the occurrence of a
manifestation unless
an additional principle
connects the two relations.

\paragraph{Proposition 14.3: Manifestation does not automatically entail Expression.}

The manifestation relation
alone does not establish
the occurrence of an
expression unless
an additional principle
connects the two relations.

\paragraph{Proposition 14.4: Expression and Manifestation may be related without being identical.}

The existence of a
specified relation between
expression and manifestation
does not establish
that the two relations
are equivalent.

\paragraph{Proposition 14.5: Expression and Manifestation are not definitions of Identity.}

Neither relation may
be substituted for Identity
without an independent
ontological argument.

\paragraph{Proposition 14.6: Plurality of Expression does not entail Identity Plurality.}

Multiple expression
instances associated
with an existent do
not, without additional
premises, establish
multiple identities.

\paragraph{Proposition 14.7: Plurality of Manifestation does not entail Identity Plurality.}

Multiple manifestation
instances associated
with an existent do
not, without additional
premises, establish
multiple identities.

\paragraph{Proposition 14.8: Formal plurality does not establish ontological independence.}

The formal distinctness
of expression or
manifestation instances
does not by itself
establish that those
instances are ontologically
independent.

The propositions must
be assessed according
to their respective
justificatory grounds.

The definitions are
stipulative formal
conventions.

The restrictions on
unwarranted implication
follow from the absence
of an established bridge
principle.

The ontological interpretation
of the relations requires
the arguments developed
in the relevant conceptual
chapters.

\subsection{Models and Countermodels}

The formal distinctions
introduced here can be
examined through models.

Let:

\[
    \mathfrak{M}
    =
    \langle
        D_E,D_X,D_M,
        \operatorname{Expr},
        \operatorname{Man},
        \operatorname{Rel}_{XM}
    \rangle.
\]

This structure represents
a model in which the
domains and relations
are assigned interpretations.

A model may satisfy
both expression and
manifestation relations
for a particular existent.

For example:

\[
    \operatorname{Expr}(a,e_1)
\]

and:

\[
    \operatorname{Man}(a,m_1).
\]

It may also satisfy:

\[
    \operatorname{Rel}_{XM}(e_1,m_1).
\]

Such a model represents
a case in which an
expression instance
and a manifestation
instance are related.

It does not establish
that every expression
must be manifested
or that every manifestation
must be expressed.

A second model may
satisfy an expression
relation without satisfying
a corresponding manifestation
relation under the specified
interpretation.

A third model may
satisfy a manifestation
relation without satisfying
a corresponding expression
relation.

These models are useful
for testing whether
the proposed implications
follow from the formal
definitions and axioms.

If a model satisfies
the premises of a proposed
implication while failing
its conclusion, the
implication is not
entailed by those
premises.

This result concerns
formal entailment.
It does not establish
that the relevant
ontological configuration
is actual or possible
in every sense.

Formal possibility
must remain distinct
from ontological possibility.

\subsection{Consistency and Independence}

The separate formal
representation of Expression
and Manifestation permits
the examination of their
logical independence.

Logical independence
must be distinguished
from ontological
independence.

Two predicates may
be independently
satisfiable in a formal
model without representing
ontologically independent
realities.

Conversely, the fact
that two concepts are
analytically distinguishable
does not by itself
establish that every
formal interpretation
must allow their
independent satisfaction.

The formal system
must therefore identify
the exact sense in
which independence
is being asserted.

At least four questions
must be distinguished:

\begin{enumerate}
    \item Can Expression
    be represented without
    a corresponding
    Manifestation relation?
    \item Can Manifestation
    be represented without
    a corresponding
    Expression relation?
    \item Are there additional
    principles that connect
    the two relations?
    \item Do those principles
    express logical,
    conceptual, or
    ontological dependence?
\end{enumerate}

The answers may differ
according to the adopted
interpretation.

No universal conclusion
about the independence
or dependence of
Expression and Manifestation
should be adopted
without specifying
the relevant level
of analysis.

\subsection{Limits of Formalization}

The formal reconstruction
provides a framework
for distinguishing
Expression and Manifestation
and for examining
their possible relations.

It does not provide
a complete ontological
account of either concept.

Several issues remain
beyond the reach of
the present formal
vocabulary:

\begin{enumerate}
    \item the complete
    ontological meaning
    of Expression;
    \item the complete
    ontological meaning
    of Manifestation;
    \item the conditions
    under which an
    expression is manifested;
    \item the conditions
    under which a
    manifestation is
    associated with
    an expression;
    \item the relation
    between expression
    and determination;
    \item the relation
    between manifestation
    and determination;
    \item the relation
    between different
    expressions of
    one existent;
    \item the relation
    between different
    manifestations of
    one existent;
    \item the conditions
    under which expression
    or manifestation
    plurality may have
    implications for
    identity;
    \item the extent
    to which the formal
    domains correspond
    to ontological
    categories.
\end{enumerate}

These issues require
further ontological
argument and cannot
be resolved merely
by increasing the
number of formal
symbols.

Formal precision
is necessary for
the reconstruction,
but formal notation
must remain accountable
to the ontological
meaning of the
concepts it represents.

\subsection{Open Problems}

The following questions
remain open.

First, what is the
most adequate positive
account of Expression
that preserves its
distinctive ontological
role?

Second, what is the
most adequate positive
account of Manifestation,
and which conditions
are essential to it?

Third, what kinds
of relation may connect
Expression and Manifestation
without collapsing
one into the other?

Fourth, can an existent
be expressed without
being manifested under
a specified interpretation?

Fifth, can an existent
be manifested without
a corresponding
expression relation?

Sixth, what is the
relation between
Expression, Manifestation,
and Determination?

Seventh, under what
conditions, if any,
may differences in
Expression or Manifestation
bear upon questions
of identity?

Eighth, how should
the ontology distinguish
formal plurality of
expression or manifestation
from ontological plurality?

Ninth, which aspects
of Expression and
Manifestation require
a formal treatment
that goes beyond
ordinary predicate
and relational logic?

These questions
must be addressed
without introducing
cognitive dependence
as a universal
condition of manifestation
or reducing Expression
to a merely linguistic
or representational
operation.

\subsection{Concluding Principle}

Expression and Manifestation
are formally distinguishable
ontological relations
whose possible connections
must be specified
rather than presumed.

Their formal reconstruction
requires separate
representations, explicit
conditions of application,
and independently
justified bridge principles
where their relations
are connected.

The governing principle
may therefore be stated
as follows:

\begin{quote}
Expression and Manifestation
may be related within
the ontological articulation
of an existent, but
neither relation may
be substituted for
the other, nor may
either be identified
with Identity or
Determination without
an independent argument.
\end{quote}

This principle preserves
the conceptual and
formal distinctions
required for the
continued reconstruction
of the ontology.

The next section
examines the formal
reconstruction of
Identity and Continuity,
with particular attention
to the conditions
under which claims
of continuity may
be formulated without
reducing identity
to persistence,
stability, or
unchanging sameness.

% =============================================================================
% END OF SECTION 14 OF CHAPTER 07
% =============================================================================